\section{Conclusion}
\label{sec:conclusion}
CreditXplain turns a credit risk model into an explainable decision process. Eligibility rules are applied before the
model, a three-band router sends uncertain cases to a loan officer, every decision is explained with SHAP and with a
local LIME that agrees with it, and applicants who are not approved receive recourse options. On the Home Credit
dataset, the model reaches a test ROC-AUC of 0.767, local LIME raises the rank correlation of LIME with SHAP from 0.02
to 0.64, and the applicant's full path through the application takes 158~ms.

The central result concerns recourse. Standard recourse options are fragile: most fail on at least one retrained
model. Requiring agreement of bootstrap-retrained models more than doubles the share of options that survive retraining,
and the hybrid used in the application gives more rejected applicants an option that is both usable and durable than
the standard search, random search or a baseline engine. A further experiment shows that removing age from the
model widens the age gap, which supports applying the age policy as explicit rules.

Future work includes a study with loan officers and applicants, validation on data with real application dates and on
other lenders' data, causal models of how the suggested changes interact, and monitoring of recourse validity and
recourse burden across groups while the system is in use.
